% ======================================================================
%  53a_dialogue_A.tex — 附录 D 第一部分《舞台与机器》问答录
%  底本：KW 2006（hep-th/0604151），§1--§5，印刷 pp.2--77
%  只收原文明确陈述的 A 级内容；每条答案可回原文核对，短引 <=120 字符
% ======================================================================

\subsection{舞台与机器：\S1--\S5（pp.2--77）}

以下二十问按原文顺序走一遍 \S1--\S5：从 $\mathcal N=4$ 超对称 Yang--Mills 的场内容与 $S$-对偶，经过扭化、以 $\PP^1$ 参数化的一族拓扑场论与典范参数 $\Psi$，到紧化给出的 Hitchin 模空间 $\MH$，最后到二维拓扑场论、generalized complex geometry，以及 $S$-对偶的镜像对称解释。答案只复述原文明确写下的内容，并附原文短引与印刷页锚点，可逐条回原文核对。原文只是暗示、没有明说的，不在此列。

\QAq{这篇论文的总目标是什么？}
\QAa{论文的目标是用四维超对称规范理论来理解复 Riemann 面的几何 Langlands 对应；原文在导言中把这套计划概括为六个主要想法，其中一个想法正是：在参数的某个特殊值上，四维 $S$-对偶表现为二维镜像对称。}
\QAs{src p.2}{``Our focus in the present paper is on the geometric Langlands program for complex Riemann surfaces.''}

\QAq{$\mathcal N=4$ 超对称 Yang--Mills 是怎么构造的？场有哪些？}
\QAa{原文从十维超对称 Yang--Mills 作维数约化构造它：十维的规范场 $A$（$G$-丛 $E$ 上的联络）与费米场 $\lambda$ 约化到四维后，$A_\mu$ 仍是规范场，而 $A_I$（$I\ge 4$）成为六个标量 $\varphi_i=A_{i+4}$。原文随后把这六个标量与一个 $G$-丛上的联络作为四维理论的场内容。}
\QAs{src p.6}{``N = 4 super Yang-Mills is most easily constructed by dimensional reduction from ten dimensions.''}

\QAq{$S$-对偶的核心断言是什么？}
\QAa{原文断言存在一个额外的量子对称性：它反演复耦合 $\tau$，把 $G$ 换成 Langlands 对偶群 $\Ldual$，并交换电与磁的电荷。对 simply-laced 的 $G$，$\tau\to-1/\tau$ 与经典的 $\tau\to\tau+1$ 生成 $SL(2,\ZZ)$；非 simply-laced 时第二生成元改为 $\tau\to-1/n_g\tau$。}
\QAs{src p.11}{``an additional quantum symmetry that inverts $\tau$, exchanges $G$ with $\Ldual$''}

\QAq{“扭化”在做什么？为什么选 GL 扭化？}
\QAa{扭化是指用 $\mathrm{Spin}(4)\times\mathrm{Spin}(6)$ 的另一个子群 $\mathrm{Spin}'(4)$ 替换原来的 $\mathrm{Spin}(4)$，使自旋表示中出现不变向量，从而留下一个不变超荷。原文指出，在 $\mathcal N=4$ 理论可用的几种扭化里，与几何 Langlands 对应相关的那一种被称为 GL 扭化。}
\QAs{src p.17}{``It turns out to be the twist relevant to the geometric Langlands program, and we will call it the GL twist.''}

\QAq{不变超荷 $Q$ 从哪里来？为什么 $Q^2=0$？}
\QAa{在 $S_+$ 中任取一个 $\mathrm{Spin}'(4)$ 不变向量，就得到一个不变超荷 $Q$；由于超对称代数满足 $\{Q,Q\}=\{\bar Q,\bar Q\}=0$，它自动满足 $Q^2=0$。原文由此把物理问题换成对 $Q$-上同调的研究。}
\QAs{src p.18}{``It will automatically obey $Q^2 = 0$, by virtue of (2.15).''}

\QAq{扭化后玻色场变成了什么？}
\QAa{原文说扭化后的玻色场是一个规范场，局域地写为取值在伴随表示中的 one-form $A=A_\mu dx^\mu$，再加上第二个伴随值 one-form $\varphi=\varphi_\mu dx^\mu$，以及复标量 $\sigma,\bar\sigma$。原来的六个标量中，$\varphi_0,\dots,\varphi_3$ 组成这个 one-form，$\sigma,\bar\sigma$ 由 $\varphi_4,\varphi_5$ 组合而成。}
\QAs{src p.19}{``a gauge field, which locally is an adjoint-valued one-form ... along with a second adjoint-valued one-form''}

\QAq{为什么会出现一族以 $\PP^1$ 参数化的拓扑场论？}
\QAa{不变超荷取 $\epsilon=u\epsilon_\ell+v\epsilon_r$，只有比值 $t=v/u$ 有意义，于是所有选择构成一个 $\PP^1$；原文说这就给出以 $\PP^1$ 参数化的一族拓扑场论。$t$ 有自然的仿射含义，且 $t$ 与 $-t$ 给出等价的拓扑场论，后者来自 $\mathrm{SU}(4)_R$ 中心的作用。}
\QAs{src p.22}{``Thus, the topological field theories with parameter $t$ or $-t$ are equivalent.''}

\QAq{这些拓扑场论的超对称方程是什么？$t=\pm i$ 时为什么变成平坦联络？}
\QAa{要求所有费米场变分为零，得到 $(F-\varphi\wedge\varphi+tD\varphi)^+=0$、$(F-\varphi\wedge\varphi-t^{-1}D\varphi)^-=0$ 与 $D*\varphi=0$。在 $t=\pm i$ 时令 $\mathcal A=A+i\varphi$，前两式等价于 $\mathcal F=0$，即 $\mathcal A$ 是平坦的 $G_{\mathbb C}$ 联络，从而给出同态 $\pi_1(M)\to G_{\mathbb C}$。}
\QAs{src p.25}{``a solution of (3.29) at $t = \pm i$ determines a complex-valued flat connection''}

\QAq{为什么说这个理论是拓扑场论？}
\QAa{原文的论证是：作用量 $I_0$ 形如 $\delta_\ell\delta_r V$，而度规不进入 $\delta_\ell,\delta_r$ 的定义，只进入 $V$ 的选取。因此作用量对度规的依赖是 $\delta_\ell\delta_r(\cdots)$ 的形式，这正是导向拓扑场论的关键性质。}
\QAs{src p.33}{``it also has the key property that leads to a topological field theory''}

\QAq{典范参数 $\Psi$ 是什么？为什么理论只依赖它？}
\QAa{原文说理论只依赖 $t$ 与 $\tau$ 的一个组合，称之为典范参数：$\Psi=\frac{\tau+\bar\tau}{2}+\frac{\tau-\bar\tau}{2}\cdot\frac{t-t^{-1}}{t+t^{-1}}$。把作用量写成 $I=\delta_tV+\frac{i\Psi}{4\pi}\int\Tr F\wedge F$ 后，$\Psi$ 以外的依赖都被 $\delta_t$ 吸收，所以只有 $\Psi$ 是重要的。}
\QAs{src p.35}{``only depends on a certain combination $\Psi$ of $t$ and $\tau$, which we will call the canonical parameter''}

\QAq{$\Psi$ 有哪些特殊值？}
\QAa{原文列出：$t=\pm 1$ 时 $\Psi=\Re\tau$；一般地 $|t|=1$ 时 $\Psi$ 为实数；$t=\pm i$ 时 $\Psi=\infty$，与 $\tau$ 无关。反过来 $t^2=-(\Psi-\bar\tau)/(\Psi-\tau)$，故任意 $\Psi$ 都能在任意弱的耦合下实现。}
\QAs{src p.36}{``If $t = \pm i$, then $\Psi = \infty$, independent of $\tau$.''}

\QAq{把四维理论紧化到 $\Sigma\times C$ 上会发生什么？}
\QAa{取 $C$ 为亏格大于 $1$ 的紧 Riemann 面、$\Sigma$ 远大于 $C$，则扭化的四维 $\mathcal N=4$ 规范理论在 $\Sigma$ 上约化为一个 sigma 模型，其映射为 $\Phi:\Sigma\to\MH$，其中 $\MH$ 是 $C$ 上 Hitchin 方程的解模空间。原文随后指出，长距离的有效物理就是目标空间为 $\MH(G,C)$ 的超对称 sigma 模型。}
\QAs{src p.38}{``reduces on $\Sigma$ to a sigma-model of maps $\Phi : \Sigma \to \MH$''}

\QAq{为什么 $\MH$ 是超凯勒流形？}
\QAa{原文把 Hitchin 方程等同于三个矩映射的消失，从而 $\MH$ 是无穷维仿射空间 $W$（所有 $(A,\varphi)$）在规范变换群下的超凯勒商。超凯勒商保持超凯勒结构，所以 $\MH$ 是超凯勒流形；原文也用这一点解释二维一侧的 $(4,4)$ 超对称。}
\QAs{src p.46}{``$\MH$ is the hyper-Kahler quotient of the infinite-dimensional space $W$ by the group of gauge transformations''}

\QAq{$\MH$ 的三个复结构 $I,J,K$ 分别看到什么？}
\QAa{$\MH$ 上有满足 $IJ=-JI=K$ 及循环置换的三个复结构 $I,J,K$，还有以 $\PP^1_h$ 参数化的一族复结构 $I_w$，其中 $J=I_{-i}$、$K=I_1$。在复结构 $I$ 下全纯数据是 Hitchin 对 $(E,\varphi)$；在复结构 $J$ 下，$\MH$ 是稳定同态 $\vartheta:\pi_1(C)\to G_{\mathbb C}$ 的模空间 $Y$（Corlette--Donaldson 定理）。}
\QAs{src p.51}{``in complex structure $J$, is the moduli space $Y$ that parametrizes stable homomorphisms''}

\QAq{Hitchin 纤维化是什么？底空间与纤维有多大？}
\QAa{把 Hitchin 对 $(E,\varphi)$ 送到 $w=\Tr\varphi^2$（一般 $G$ 改用独立的 Casimir 不变量），得到映到复向量空间 $B$ 的 Hitchin 纤维化 $\pi:\MH\to B$。$B$ 的复维是 $(g-1)\dim G$，等于 $\MH$ 复维的一半；纤维 $F$ 的维数与 $B$ 相同，且 Hitchin 纤维化是满射。}
\QAs{src p.53}{``The Hitchin fibration, as it is most commonly called, is the map $\pi : \MH \to B$''}

\QAq{二维拓扑场论什么时候是 A 模型、什么时候是 B 模型？}
\QAa{二维拓扑场论由一对复结构 $(J_+,J_-)$ 决定：$J_+=J_-$ 时是 B 模型，$J_+=-J_-$ 时是 A 模型，其余情形可用 generalized complex geometry 化到 A 模型。从四维约化看 $(w_+,w_-)=(-t,t^{-1})$，所以 $t=\pm i$ 恰给出 B 模型，$t$ 为实数时给出 A 模型。}
\QAs{src p.61}{``For a B-model, we want $w_+ = w_-$. So the condition is $t^{-1} = -t$, which occurs precisely for $t = \pm i$.''}

\QAq{generalized complex geometry 在这里起什么作用？}
\QAa{原文用它把二维拓扑场论写成广义复结构：由 $\MH$ 的超凯勒度量与一对点 $J_+,J_-\in\PP^1_h$ 得到一对广义复结构，其中决定模型的那个只依赖 $J$、$\omega_K$ 与 $\omega_I$，这三个量与 $C$ 的复结构无关，故约化得到的这族理论不依赖 $C$ 的度规。进一步，一般情形下它是 A 模型再加一个 $B$ 场变换，由此给出 $\Psi$ 的二维解释。}
\QAs{src p.67}{``The first important observation is that $I$, $K$, and $\omega_J$ have disappeared.''}

\QAq{$S$ 交换 $\Psi=0$ 与 $\Psi=\infty$；这两个端点各是什么模型？}
\QAa{原文取 $\tau$ 为纯虚数：$\Psi=\infty$（$t=i$）给出复结构 $J$ 的 B 模型，与 $\tau$ 的取值无关；$S$ 把 $t$ 送到 $1$，得到复结构 $K$ 的 A 模型，即 $\Psi=0$。因此 $S$ 把 $\MH(G,C)$ 的 B 模型换成 $\MH(\Ldual,C)$ 的 A 模型。}
\QAs{src p.70}{``that maps $\Psi \to -1/n_g \Psi$ and hence exchanges $\Psi = 0$ and $\Psi = \infty$''}

\QAq{$S$-对偶怎样在 Hitchin 纤维化上实现镜像对称？}
\QAa{原文要证的是 $S$-对偶在 Hitchin 纤维化 $\pi:\MH\to B$ 的底空间上经典地作用，而在纤维上作 $T$-对偶。零膜 $B_p$ 是 (B,B,B) 型，$S$ 把它变成支撑在某个纤维 $\tilde F$ 上的 (B,A,A) 型膜，两侧模空间分别是 $F$ 与 $\tilde F$ 上秩一平坦 Chan--Paton 丛的复环面，$S$ 在它们之间建立同构。原文总结说这正是 SYZ 的环面纤维对偶，只是底空间之间的映射 $\Xi:B\to{}^LB$ 在 SYZ 中通常被当作平凡。}
\QAs{src p.71}{``$S$-duality acts classically on the base of the Hitchin fibration''}

\QAq{$\MH$ 上有哪些类型的膜？$S$-对偶怎样交换它们？}
\QAa{原文列出 (B,B,B) 型膜，即在 $I,J,K$ 三个复结构下都是 B 型，以及 (B,A,A)、(A,B,A)、(A,A,B) 型膜，即在一个复结构下是 B 型、在另外两个下是 A 型。$S$ 把 (B,B,B) 型换成 (B,A,A) 型，而 (A,B,A) 与 (A,A,B) 型各自换成同一类型的膜。}
\QAs{src p.77}{``(B, B, B)-branes, which are B-branes in each complex structure $I$, $J$, $K$''}
